% Figure: the three boundary-layer length scales read off one velocity profile.
% The streamwise velocity u/U_inf rises from zero at the wall to the freestream
% value, reaching 99 percent at y = delta. The displacement thickness delta*
% is the area of the velocity deficit (1 - u/U_inf); the momentum thickness
% theta weights that deficit by u/U_inf. The sketch uses the cubic profile
% u/U = 1.5(y/delta) - 0.5(y/delta)^3 so the shape is the one worked example 10
% integrates.
\begin{figure}[htbp]
\centering
\begin{tikzpicture}
\begin{axis}[
  width=11cm, height=7cm,
  xlabel={velocity ratio $u/U_\infty$},
  ylabel={wall distance $y/\delta$},
  xmin=0, xmax=1.15, ymin=0, ymax=1.25,
  legend pos=south east,
  legend cell align=left,
  grid=both, grid style={enggray!18},
  minor grid style={enggray!8},
  tick label style={font=\scriptsize},
  label style={font=\small},
  legend style={font=\scriptsize},
]
% Freestream reference line at u/U = 1.
\addplot[enggray, dashed, thick] coordinates {(1,0) (1,1.25)};
\node[font=\scriptsize, enggray, anchor=south] at (axis cs:1,1.2) {$U_\infty$};
% Cubic velocity profile, plotted as u/U on x and y/delta on y.
% Sample y/delta = t from 0 to 1, u/U = 1.5 t - 0.5 t^3; above delta, u = U.
\addplot[engblue, very thick, samples=60, domain=0:1, variable=\t]
  ({1.5*\t - 0.5*\t^3},{\t});
\addplot[engblue, very thick] coordinates {(1,1) (1,1.25)};
\addlegendentry{velocity profile $u/U_\infty$}
% delta marker at y/delta = 1.
\draw[enggray, thick, {Latex}-] (axis cs:0.05,1) -- (axis cs:0.25,1);
\node[font=\scriptsize, anchor=west] at (axis cs:0.26,1) {$y=\delta$ ($u=0.99U_\infty$)};
% displacement-thickness annotation arrow inside the deficit region.
\draw[engred, thick, -{Latex}] (axis cs:0.72,0.55) -- (axis cs:0.93,0.62);
\node[font=\scriptsize, engred, anchor=east, align=right]
  at (axis cs:0.70,0.52) {velocity deficit\\ $1-u/U_\infty$\\ (sets $\delta^*$)};
% momentum-thickness annotation near the wall where u(U-u) is largest.
\draw[engorange, thick, -{Latex}] (axis cs:0.35,0.20) -- (axis cs:0.52,0.30);
\node[font=\scriptsize, engorange, anchor=west, align=left]
  at (axis cs:0.53,0.30) {momentum flux deficit\\ $\tfrac{u}{U_\infty}(1-\tfrac{u}{U_\infty})$ (sets $\theta$)};
\end{axis}
\end{tikzpicture}
\caption[Boundary-layer thicknesses from one profile]{One velocity profile
carries three length scales. The boundary-layer thickness $\delta$ is the
height at which $u$ reaches $99\%$ of the freestream. The displacement
thickness $\delta^*$ is the integral of the velocity deficit $1-u/U_\infty$,
the distance the wall would have to move out to pass the same mass as an
inviscid flow. The momentum thickness $\theta$ is the integral of
$(u/U_\infty)(1-u/U_\infty)$, which weights the deficit by the local velocity
and measures the momentum the layer has lost to the wall. For the cubic profile
drawn here, $\delta^*=\tfrac38\delta$ and $\theta=\tfrac{39}{280}\delta$, so the
three scales stand in a fixed ratio fixed by the profile shape.}
\label{fig:vol03ch12:thicknesses}
\end{figure}
